% ================================================================
% NetFix AI — Root Cause Analysis chapter
% Drop-in chapter for the shared thesis main.tex / thesis_setup.tex.
% No document class or package imports are intentionally included here.
% ================================================================

\chapter{Root Cause Analysis and Diagnostic Intelligence}
\label[chapter]{chap-rca}

% Local JSON language definition. The listings package is loaded globally
% by thesis_setup.tex.
\lstdefinelanguage{json}{
  basicstyle=\ttfamily\scriptsize,
  numbers=left,
  numberstyle=\tiny,
  stepnumber=1,
  numbersep=7pt,
  showstringspaces=false,
  breaklines=true,
  frame=single,
  columns=fullflexible
}

\section{System Overview and Architecture}

This chapter details the reasoning and explanation layer of the Drive-Test QoE Anomaly Detection \& Root Cause Analysis (RCA) system. While the upstream anomaly-detection pipeline converts raw, second-by-second LTE drive-test observations into consolidated engineer-facing incidents, the RCA layer provides the engineering intelligence that transforms those incidents into actionable, explainable, and auditable network diagnoses.

The core design philosophy rests on two strict operational principles:

\begin{enumerate}
\item \textbf{KPI as Evidence, Not Cause:} Low RSRP or high Resource Block (RB) utilization are physical measurements, not root causes. The system maps measured evidence to explicit physical mechanisms such as \textit{Weak Serving Cell} or \textit{Cell Congestion}.

\item \textbf{Deterministic Rules Over Generative Hallucination:} The root-cause decision is made by a priority-ordered, deterministic rule engine grounded in LTE engineering knowledge and 3GPP-defined radio procedures. The Large Language Model (LLM) is positioned strictly downstream as a diagnostic narrative layer; the deterministic cause is finalized before prompting, and the model is not authorized to alter the assigned cause.
\end{enumerate}

The LTE procedures and KPI semantics used by the rule base are consistent with the E-UTRA RRC, physical-layer, and MAC specifications \cite{3gpp36331,3gpp36213,3gpp36321}. The numerical diagnosis gates used in this project are engineering rules calibrated for the project workflow; they should not be interpreted as universal 3GPP optimization thresholds.

\begin{infobox}{RCA processing chain}
\centering
\textbf{Drive-Test KPIs}
$\rightarrow$
\textbf{Upstream Anomaly Detection}
$\rightarrow$
\textbf{Incident Consolidation}
$\rightarrow$
\textbf{Deterministic RCA Engine}
$\rightarrow$
\textbf{Solution Mapping}
$\rightarrow$
\textbf{LLM Diagnostic Narrative}
\end{infobox}

\subsection{Scope and Population Metadata}

\begin{table}[htbp]
\centering
\small
\begin{tabularx}{\textwidth}{>{\bfseries}p{4.0cm} p{3.2cm} X}
\toprule
Metric / Parameter & Value & Description\\
\midrule
Primary Problem Scope & LTE DL Throughput Degradation & Focus on QoE degradation\\
Analysis Unit & Incident-Level Aggregate & Aggregated anomaly episodes with median KPI state\\
Total Incident Population & 1,382 & Analyzed incident cases\\
Underlying Drive-Test Rows & 2,824 & Raw observation samples represented\\
Unique Serving Cells & 423 & Distinct serving-cell identities represented in the RCA handoff\\
Unique Test Sessions & 254 & Distinct drive-test runs\\
Observation Window & Aug. 26, 2025 (23:37) $\rightarrow$ Aug. 29, 2025 (00:36) & Approximately 3-day monitoring period\\
Decision Layer & Deterministic, Priority-Ordered, KPI-Driven & Knowledge base via \texttt{causes.json} and \texttt{kpi\_rules.json}\\
Explanation Layer & LLM-Assisted Narrative Generation & Converts structured JSON verdict into engineer reports\\
\bottomrule
\end{tabularx}
\caption{Scope and population metadata for the complete RCA evaluation.}
\label[tab]{rca-scope}
\end{table}

Figure~\ref{rca-dataset-overview} provides a compact descriptive view of the incident-level RCA data. The deterministic evaluation population used throughout this chapter contains \textbf{1,382 incidents} and \textbf{2,824 underlying drive-test rows}, as summarized in Table~\ref{rca-scope}.

\begin{figure}[htbp]
\centering
\includegraphics[width=0.98\linewidth]{RCA_01_dataset_overview.png}
\caption{Incident-level RCA dataset overview showing final population metadata, severity distribution, dominant RCA classifications, and median KPI context.}
\label[fig]{rca-dataset-overview}
\end{figure}

\FloatBarrier

\section{Contributor Responsibilities and Boundary Definition}

The RCA work package covers the complete pipeline between the upstream anomaly-detection handoff and the final downstream diagnostic representation.

\begin{table}[htbp]
\centering
\small
\begin{tabularx}{\textwidth}{>{\bfseries}p{3.4cm} p{6.0cm} X}
\toprule
Functional Module & Technical Responsibility & Primary Technical Output\\
\midrule
Anomaly Monitoring & Consolidates raw anomaly rows into persistent, repeated incident cases. & Prioritized Incident Cases\\
RCA Cause Taxonomy & Defines the hierarchical priority structure of engineering causes. & Priority Cause Tree\\
Knowledge Base Architecture & Designs the decoupled three-JSON system (\texttt{causes}, \texttt{kpi\_rules}, \texttt{solutions}). & Maintainable Knowledge Base\\
KPI Reasoning Engine & Implements rule evaluation logic (\texttt{MATCH\_ALL}/\texttt{MATCH\_ANY}) against evidence. & Primary Cause ID + Evidence Trace\\
Evidence Extraction & Extracts and formats exact KPI values supporting the selected cause. & \texttt{key\_kpi\_evidence} Payload\\
Remediation Mapping & Links confirmed cause IDs to standardized telecom remediation actions. & \texttt{recommended\_solution} Entry\\
LLM Diagnostic Layer & Prompts the LLM to render structured RCA JSON into natural diagnostic prose. & Human-Readable Diagnostic Narrative\\
\bottomrule
\end{tabularx}
\caption{RCA responsibilities and technical outputs.}
\label[tab]{rca-responsibilities}
\end{table}

\textbf{Scope Boundary:} Upstream anomaly detection is treated as an input. This chapter explains \textit{why} degradation occurred, \textit{what measured evidence} supports the verdict, and \textit{what actions} engineers should take. Planning Intelligence remains an independent data stream and is never used as proof for drive-test anomalies.

\section{Data and KPI Taxonomy}

The RCA engine evaluates multi-dimensional drive-test telemetry organized into distinct functional KPI categories.

\begin{table}[htbp]
\centering
\scriptsize
\begin{tabularx}{\textwidth}{>{\bfseries}p{3.0cm} p{5.0cm} X}
\toprule
KPI Group & Specific Metrics / Indicators & Engineering Significance in RCA\\
\midrule
Throughput (Symptom) & Actual DL Throughput, Expected $p_{75}$ Throughput, Throughput Gap & Defines performance drop severity and baseline gap\\
Coverage & RSRP (Serving Cell) & Identifies weak coverage and cell-edge conditions\\
Radio Quality & RSRQ, SINR & Separates inter-cell interference from signal attenuation\\
Link Adaptation & CQI, MCS, BLER & Evaluates channel state reporting and modulation efficiency\\
Resource Allocation & RB Utilization, Cell PRB Load, Throughput/RB & Detects air-interface congestion versus scheduler under-grant\\
Carrier Aggregation (CA) & CA Active Flag, Carrier Count, SCell1 RSRP/SINR/CQI/BLER & Identifies secondary carrier underperformance or deactivation\\
Mobility & HO Count ($\pm5$ s), Dwell Time, Serving-Neighbor RSRP Delta, UE Speed & Evaluates handover timing, ping-ponging, and neighbor relations\\
Access & RACH Failure Flag, RACH Attempts, Access Delay & Identifies random-access channel contention and link setup failures\\
MIMO / Spatial & Rank Indicator (RI), Layer Utilization & Detects rank fallback (RI=1) under strong SINR conditions\\
\bottomrule
\end{tabularx}
\caption{KPI taxonomy used by the RCA engine.}
\label[tab]{rca-kpi-taxonomy}
\end{table}

Figure~\ref{rca-kpi-distributions} gives the empirical distribution of the principal radio and link-adaptation KPIs in the incident dataset. The broad distributions and long tails illustrate why RCA is based on combinations of evidence rather than a single degraded KPI.

\begin{figure}[htbp]
\centering
\includegraphics[width=0.98\linewidth]{RCA_02_core_kpi_distributions.png}
\caption{Core radio KPI distributions across the incident dataset: RSRP, RSRQ, SINR, CQI, MCS, and BLER.}
\label[fig]{rca-kpi-distributions}
\end{figure}

\FloatBarrier

\section{RCA Reasoning Framework and Taxonomy}

To prevent a single degraded KPI, such as high BLER, from falsely triggering a diagnosis, the framework enforces a \textbf{priority-ordered evaluation hierarchy}. Causes higher in the hierarchy are evaluated first; when their defining gate conditions are satisfied, the engine can stop at the highest-priority confirmed mechanism, preserving deterministic and traceable reasoning.

Figure~\ref{rca-cause-taxonomy} summarizes the engineering cause tree around the central symptom of throughput degradation. It separates resource limitation, coverage and radio quality, cell health, mobility and access, MIMO performance, and carrier aggregation into explicit cause families supported by measurable KPI evidence.

\begin{figure}[htbp]
\centering
\includegraphics[width=0.99\linewidth]{RCA_03_throughput_cause_taxonomy.png}
\caption{Engineering RCA cause taxonomy for LTE throughput degradation and the supporting KPI/evidence branches.}
\label[fig]{rca-cause-taxonomy}
\end{figure}

\subsection{Priority 1: Resource Limitation}

\textbf{Cell Congestion (C01):} Identifies capacity-limited cases under confirmed download demand, where high resource use and a substantial throughput shortfall indicate a genuine resource ceiling rather than a radio-quality limitation. Additional evidence can identify high cell load, low rank, unavailable CA, or low operational bandwidth.

\textbf{Poor Spectral Efficiency (C02):} Identifies scheduler under-grant when demand is present but the available radio link is not being translated into sufficient resource allocation. Low MCS, low SINR, and high BLER are retained as supporting evidence when applicable.

\subsection{Priority 2: Coverage and Radio Quality}

\textbf{Weak Serving Cell (C03):} Identifies poor serving-cell coverage through the configured weak-RSRP gate. Low RSRQ, low MCS, high BLER, and cell-edge context are retained as supporting evidence rather than treated as independent root causes.

\textbf{Interference (C04):} Identifies degraded radio quality despite adequate serving-signal power. Degraded SINR is the principal interference signature, while CQI, RSRQ, and BLER provide corroborating evidence when available.

\subsection{Priority 3: Outage / eNodeB Fault}

\textbf{Serving Cell / eNodeB Fault (C05):} Reserved for a declared Radio Link Failure while the UE's own radio conditions remain normal, indicating a likely serving-cell or hardware-side fault rather than coverage or interference. Sub-causes that require OSS, RRC-release, power, or neighbor-relation data remain non-computable when those sources are unavailable.

\textbf{Outage (C06):} Identifies a throughput or session collapse while the UE's own radio conditions remain normal. This branch deliberately represents a general network-side service disruption: drive-test evidence alone does not determine whether the serving site, a neighboring site, transport, or another network element is at fault. Supporting evidence such as a serving-cell utilization increase or a neighbor-availability drop is retained as an investigation lead rather than treated as independent confirmation.

\subsection{Priority 4: Mobility}

\textbf{Handover Problem (C08):} Captures handover-adjacent degradation through supporting patterns such as ping-pong handover, recurring wrong-neighbor behavior, late handover, and high-speed mobility context.

\textbf{Access Problem (C09):} Captures Random Access Channel (RACH) failure or contention signatures. More specific explanations that require uplink interference or RACH transmit-power telemetry remain non-computable when those measurements are unavailable.

\subsection{Priority 5: MIMO}

\textbf{MIMO Rank Fallback (C10):} Identifies poor spatial multiplexing when link quality is strong but the Rank Indicator falls back to Rank 1, pointing toward antenna-path, correlation, MIMO-capability, or adaptation issues.

\subsection{Priority 6: Carrier Aggregation}

\textbf{Weak Secondary Cell (C11):} Identifies active Carrier Aggregation whose secondary carrier is limited by poor SCell radio quality. Low SCell SINR, low SCell RSRP, and low or zero SCell RB allocation are retained as supporting evidence.

\textbf{SCell Release / Deactivation (C12):} Identifies secondary-carrier deactivation during an active session. High SCell BLER or low SCell CQI before the release can support a radio-quality-driven explanation.

\textbf{CA Inactive / Disabled (C13):} Identifies sessions in which Carrier Aggregation never becomes active and the UE remains on a single carrier, leaving no aggregated bandwidth or SCell radio evidence available.

\subsection{Priority 7: Fallback / Inconclusive}

\textbf{Radio Limited (C14):} A lower-priority fallback for degraded link adaptation, such as elevated BLER or low CQI, when no higher-priority cause has already been confirmed.

\textbf{Inconclusive (C15):} The definitive last-resort classification when active demand is not confirmed or no condition in the preceding hierarchy matches. It is retained for auditability but is not interpreted as a confirmed network root cause.

\FloatBarrier

\section{Knowledge Base Architecture: The Three-JSON System}

To prevent domain knowledge from being hardcoded inside Python scripts, the RCA architecture separates concerns across three JSON files linked by stable identifiers (\texttt{Cause ID} and \texttt{Solution ID}). Figure~\ref{rca-three-json} shows how the taxonomy, KPI evaluation logic, and remediation knowledge remain independently maintainable while sharing deterministic IDs.

\begin{figure}[htbp]
\centering
\includegraphics[width=0.98\linewidth]{RCA_04_three_json_architecture.png}
\caption{Three-JSON RCA knowledge-base architecture: \texttt{causes.json} defines what is evaluated and in which order, \texttt{kpi\_rules.json} defines how evidence is tested, and \texttt{solutions.json} defines what action follows a confirmed cause.}
\label[fig]{rca-three-json}
\end{figure}

\begin{table}[htbp]
\centering
\small
\begin{tabularx}{\textwidth}{>{\ttfamily}p{3.0cm} >{\bfseries}p{4.0cm} X}
\toprule
JSON File & Core System Role & Key Content and Schema\\
\midrule
causes.json & Taxonomy \& Hierarchy & Priority categories, parent cause IDs (C01--C15, excluding C07), supporting-evidence IDs, evaluation modes, reason templates (\texttt{reason\_template}), and computability notes.\\
kpi\_rules.json & Evaluation Logic \& Thresholds & gate conditions,
supporting evidence, alternatives, and structured
\texttt{\{field, op, value\}} condition trees.\\
solutions.json & Remediation Knowledge & Entries keyed directly by parent cause ID, with \texttt{recommended\_actions}, optional evidence-specific \texttt{applies\_if} conditions, and \texttt{standards\_and\_vendor\_references}.\\
\bottomrule
\end{tabularx}
\caption{Knowledge-base component summary.}
\label[tab]{rca-json-components}
\end{table}

\FloatBarrier

\section{Deterministic RCA Engine Execution Flow}

The classification script (\texttt{classify\_incidents.py}) executes a rule-driven pipeline at runtime:

\begin{enumerate}
\item \textbf{Load Artifacts:} Parses \texttt{causes.json}, \texttt{kpi\_rules.json}, \texttt{solutions.json}, and input incident records.
\item \textbf{Extract Telemetry:} Extracts raw and derived KPI fields using a safe extractor (\texttt{extract\_fields()}).
\item \textbf{Evaluate Hierarchy:} Sequentially tests cause gates from Priority 1 to Priority 7.
\item \textbf{Apply Boolean Logic:} Evaluates each parent gate using its configured \texttt{MATCH\_ALL} or \texttt{MATCH\_ANY} logic. When a parent gate fires, its supporting-evidence IDs are evaluated independently and collected before the hierarchy stops at the selected root cause.
\item \textbf{Format Reason Template:} Populates the cause's \texttt{reason\_template} with actual KPI values rounded to two decimal places.
\item \textbf{Attach Solution:} Fetches corresponding actions and standards citations from \texttt{solutions.json}.
\item \textbf{Emit Structured Output:} Writes the final JSON payload.
\end{enumerate}

\subsection{Real RCA Output Example}

To avoid mixing values from different incidents, the thesis uses a real record from the final RCA output rather than a synthetic example. The listing below reproduces one complete incident as emitted by the RCA pipeline; its diagnosis, KPI values, and recommendations all belong to the same case.

\begin{lstlisting}[language=json,caption={Real RCA incident output for INC-000001.}]
{
  "incident_id": "INC-000001",
  "diagnosis": {
    "category": "Outage / eNodeB Fault",
    "problem_name": "Outage",
    "reason": "A throughput/session collapse occurred while this UE's own radio conditions (RSRP -104.58 dBm, SINR 8.44 dB) remained normal -- pointing to a network-side service disruption rather than a radio-quality issue.",
    "cause_id": "C06",
    "severity": "Medium"
  },
  "key_kpi_evidence": {
    "rlf_event_present": false,
    "actual_throughput_mbps": 0.03,
    "expected_throughput_mbps": 0.47,
    "rsrp_dbm": -104.58,
    "rsrq_db": -9.48,
    "sinr_db": 8.44,
    "cqi": 8.0,
    "mcs": 10.0,
    "bler_pct": 11.76,
    "rb_usage_pct": 11.0,
    "cell_prb_utilization_pct": 11.0,
    "carrier_count": 1.0,
    "lte_rank": 1.0,
    "near_handover_mobility_flag": false,
    "ue_speed_kmh": 0.0,
    "serving_minus_best_neighbor_rsrp": 9.79,
    "handover_count_pm5s": 0.0,
    "near_rach_failure_flag": false,
    "neighbor_cell_availability_pct": 100.0
  },
  "recommended_solution": {
    "recommended_actions": [
      "Check OSS alarms/counters for both the serving site AND its immediate neighbors during this exact time window -- this pattern (throughput/session collapse with no radio explanation) cannot on its own tell you which site is at fault.",
      "Cross-reference against site maintenance, reset, or transport-link logs for the same window before dispatching an RF team.",
      "No RF tuning action is appropriate here until an outage/alarm event is confirmed or ruled out on either site."
    ],
    "standards_and_vendor_references": [
      {
        "source": "3GPP TS 32.111 (Fault Management)",
        "note": "Standard practice is to correlate a throughput/traffic collapse signature against OSS alarms before pursuing an RF explanation -- and to check neighboring sites, not just the serving cell, before attributing fault."
      }
    ]
  }
}
\end{lstlisting}

The fault-management reference stored in the incident is interpreted conservatively in this thesis. 3GPP TS 32.111 is a multipart Fault Management family; the Alarm Integration Reference Point information service is specified in TS 32.111-2 \cite{3gpp32111-2}. The recommendation to inspect both the serving site and nearby sites is an engineering investigation step used by this project, not a verbatim normative requirement of TS 32.111-2. The emitted \texttt{cause\_id} and category labels are reproduced from the same final RCA artifact and are not re-derived inside the thesis.

\section{Incident-Level Empirical Results (1,382 Incidents)}

The RCA engine was executed against the dataset of 1,382 incident-level aggregates.

\subsection{RCA Classification Distribution Summary}

\begin{table}[htbp]
\centering
\tiny
\begin{tabularx}{\textwidth}{c X r r X}
\toprule
Cause ID & Engineering Cause Name & Count & Share & Dominant Category\\
\midrule
C01 & Cell Congestion (Capacity-Limited) & 365 & 26.4\% & Resource Limitation\\
C03 & Weak Serving Cell & 237 & 17.1\% & Coverage \& Radio Quality\\
C15 & Inconclusive & 236 & 17.1\% & Fallback / Inconclusive\\
C14 & Radio Limited (Link Adaptation) & 181 & 13.1\% & Fallback / Inconclusive\\
C06 & Outage & 98 & 7.1\% & Outage / eNodeB Fault\\
C13 & CA Inactive / Disabled & 92 & 6.7\% & Carrier Aggregation\\
C08 & Handover Problem & 44 & 3.2\% & Mobility\\
C11 & Weak Secondary Cell (SCell) & 40 & 2.9\% & Carrier Aggregation\\
C10 & MIMO Rank Fallback & 33 & 2.4\% & MIMO\\
C04 & Interference / Radio Link Failure & 28 & 2.0\% & Coverage \& Radio Quality\\
C09 & Access Problem (RACH Failures) & 18 & 1.3\% & Mobility\\
C05 & Serving Cell / eNodeB Fault & 5 & 0.4\% & Outage / eNodeB Fault\\
C02 & Poor Spectral Efficiency & 5 & 0.4\% & Resource Limitation\\
\midrule
\textbf{Total} & \textbf{All Categories} & \textbf{1,382} & \textbf{100.0\%} &\\
\bottomrule
\end{tabularx}
\caption{RCA classification distribution across 1,382 incidents.}
\label[tab]{rca-cause-distribution}
\end{table}

Figure~\ref{rca-severity-composition} complements the count table by showing how severity is distributed within each engineering cause. This makes it easier to distinguish frequently occurring medium-severity mechanisms from the smaller number of high-severity cases.

\begin{figure}[htbp]
\centering
\includegraphics[width=0.86\linewidth]{RCA_05_severity_by_root_cause.png}
\caption{Severity composition per RCA classification in the deterministic results.}
\label[fig]{rca-severity-composition}
\end{figure}

\subsection{Key Engineering Findings}

\begin{enumerate}
\item \textbf{Resource Saturation Dominance:} Cell Congestion (C01) represents the single largest cause (26.4\%, 365 incidents). These cases occurred under verified download demand where RBs were saturated, indicating that capacity expansion or traffic offloading is required rather than RF tuning.

\item \textbf{Coverage Limits:} Weak Serving Cell (C03) represents 17.1\% (237 incidents), confirming that RF attenuation remains a major bottleneck in the drive-test area.

\item \textbf{Operational Trust via Inconclusive Classification:} Inconclusive (C15) accounts for 17.1\% (236 incident classifications). When active demand is not confirmed or no higher-priority rule matches, the engine retains an explicit fallback result rather than forcing a network diagnosis. C15 is therefore not interpreted as a confirmed network root cause.
\end{enumerate}

\FloatBarrier

\section{Remediation Mapping Engine}

Once a root cause is established, the engine retrieves corresponding remediation guidelines from \texttt{solutions.json}.

\begin{table}[htbp]
\centering
\scriptsize
\begin{tabularx}{\textwidth}{>{\bfseries}p{3.0cm} p{4.0cm} X}
\toprule
Cause Category & Primary Engineering Root Cause & Direct Remediation Action\\
\midrule
Resource Limitation & Cell Congestion (C01) & Run a capacity audit, review scheduler fairness/weights, and consider Mobility Load Balancing toward underutilized neighbors; prioritize expansion when the cell is genuinely capacity-exhausted.\\
Coverage \& Radio & Weak Serving Cell (C03) & Audit serving-site tilt, azimuth, and transmit power; consider a new site, small cell, or repeater when a repeated coverage hole is confirmed.\\
Coverage \& Radio & Interference (C04) & Perform an RF/interference investigation and first verify that repeated measurements show adequate RSRP, preventing a weak-coverage case from being treated as interference.\\
Outage / eNodeB Fault & Serving Cell Fault (C05) & Check serving-site OSS alarms/counters and relevant mobility-failure counters for the event window; avoid RF tuning until an outage or fault is confirmed or ruled out.\\
Outage / eNodeB Fault & Outage (C06) & Correlate the event with OSS alarms/counters for the serving site and immediate neighbors, then inspect maintenance, reset, and transport-link logs before attributing the failure to a specific site.\\
Mobility & Handover Problem (C08) & Tune Event A3 hysteresis/offset and Time-to-Trigger according to the matched handover pattern; use an ANR audit when the problem recurs across multiple episodes.\\
Mobility & Access Problem (C09) & Review RACH configuration, including preamble and power-ramping parameters, and investigate uplink coverage/interference when access delay is elevated.\\
MIMO & Rank Fallback (C10) & Inspect antenna/RF paths, cross-polarization, spacing, and correlation at affected sites; check UE MIMO capability/reporting when the behavior is UE-specific.\\
Carrier Aggregation & Weak Secondary Cell (C11) & Audit RF coverage on the secondary carrier and inspect SCell scheduler weighting when allocation remains very low; do not disable CA as a substitute for correcting the SCell limitation.\\
Carrier Aggregation & CA Inactive / Disabled (C13) & Confirm UE CA capability and secondary-carrier availability, then check CA feature activation/licensing and UE-context capability flags.\\
\bottomrule
\end{tabularx}
\caption{Representative cause-to-remediation mapping derived from \texttt{solutions.json}.}
\label[tab]{rca-remediation}
\end{table}

\section{Downstream LLM Diagnostic Narration}

The Large Language Model (LLM) operates strictly as a text-generation layer downstream of the deterministic RCA engine. It translates structured JSON verdicts into concise diagnostic reports for Network Operations Center (NOC) engineers.

\subsection{Recommended Prompting Structure}

\begin{enumerate}
\item \textbf{Context \& Symptom:} State the incident ID, timestamp, serving cell, and throughput drop.
\item \textbf{Deterministic Verdict:} State the primary cause name and cause ID (C01--C15, with C07 unused in the current taxonomy).
\item \textbf{KPI Evidence:} Cite the exact numerical KPI values that triggered the decision rule.
\item \textbf{Reasoning Explanation:} Explain why the KPI combination supports this physical mechanism.
\item \textbf{Action Plan:} Present the recommended remediation steps from \texttt{solutions.json}.
\end{enumerate}

\begin{warningbox}{System Guardrail}
The LLM prompt strictly instructs the model to explain the deterministic verdict. It is forbidden from altering the assigned \texttt{cause\_id}, inventing unmeasured KPIs, or proposing remedies outside \texttt{solutions.json}.
\end{warningbox}

\section{System Validation and Explainability Audit}

The system's architecture directly addresses key engineering auditability requirements.

\begin{table}[htbp]
\centering
\small
\begin{tabularx}{\textwidth}{>{\bfseries}p{5.0cm} X}
\toprule
Engineering Question & System Audit Mechanism\\
\midrule
Why was this cause selected? & Traced to a specific priority rule gate in \texttt{kpi\_rules.json} that was satisfied.\\
Which KPIs support the decision? & All matching KPI values are explicitly stored in the \texttt{key\_kpi\_evidence} JSON payload.\\
Can thresholds be updated? & Yes. Modifying values in \texttt{kpi\_rules.json} changes behavior without editing Python code.\\
Can remedies be customized? & Yes. Updating \texttt{solutions.json} updates engineering recommendations system-wide.\\
How are data gaps handled? & Missing fields are flagged as \texttt{NOT\_COMPUTABLE}, preventing false conclusions.\\
Can an engineer override the output? & Yes. The structured JSON retains full evaluation traces (\texttt{evaluation\_trace}) for manual review.\\
Can the LLM alter the deterministic cause? & The cause decision is finalized before the LLM prompt is constructed, and the LLM is instructed not to alter the assigned cause.\\
\bottomrule
\end{tabularx}
\caption{Explainability and auditability mechanisms.}
\label[tab]{rca-auditability}
\end{table}

\section{Future Work and Enhancements}

\begin{enumerate}
\item \textbf{Row-Level Sequential RCA:} Extend the engine from incident-level aggregates to row-level time-series data to capture real-time state transitions, such as mid-session SCell drops.

\item \textbf{OSS \& Alarm Integration:} Ingest eNodeB hardware alarms, MME paging logs, and backhaul status counters to resolve C05/C06 outage sub-causes.

\item \textbf{Dynamic Cell-Specific Baselines:} Replace selected global engineering thresholds with cell-specific historical percentiles ($p_{10}/p_{50}$) once sufficient historical data is available.

\item \textbf{Enhanced Evidence Scoring:} Implement a multi-factor confidence score based on the number and independence of supporting KPI indicators.

\item \textbf{LLM Verification Guard:} Implement an automated validation check that compares generated text against the source JSON to ensure factual alignment.
\end{enumerate}

\section{RCA Chapter Conclusion}

This contribution establishes an auditable, rule-driven Root Cause Analysis framework for LTE drive-test throughput anomalies. By enforcing a deterministic decision layer backed by a modular three-JSON knowledge base, the system converts measured telemetry into clear engineering diagnoses linked to actionable remediation steps.

Applied to the 1,382-incident evaluation dataset, the engine identified \textbf{Cell Congestion} (26.4\%) and \textbf{Weak Serving Cell} (17.1\%) as the primary network bottlenecks, while maintaining operational trust by retaining \textbf{Inconclusive} fallback classifications (17.1\%) instead of forcing unsupported network diagnoses. The downstream LLM layer converts these structured verdicts into clear diagnostic prose without sacrificing deterministic rigor.

\begin{successbox}{Technical Pipeline Summary}
\centering
\texttt{Detected Anomaly}
$\rightarrow$
\texttt{Monitored Incident}
$\rightarrow$
\texttt{KPI Evidence Extraction}
$\rightarrow$
\texttt{Deterministic Rule Match}
$\rightarrow$
\texttt{Remediation Retrieval}
$\rightarrow$
\texttt{LLM Diagnostic Narrative}
\end{successbox}

\section{Suggested Schema for Downstream Integration}

\begin{table}[htbp]
\centering
\small
\begin{tabularx}{\textwidth}{>{\ttfamily}p{4.35cm} >{\ttfamily}p{2.65cm} X}
\toprule
Field Name & Data Type & Purpose\\
\midrule
incident\_id & String & Unique source incident identifier\\
timestamp & ISO-8601 String & Incident start/end window\\
serving\_cell\_eci & Float/String & E-UTRAN Cell Identifier (ECI)\\
primary\_cause\_id & String & Stable cause identifier (e.g., C01, C03)\\
primary\_cause\_name & String & Human-readable engineering cause name\\
severity & String & Incident severity level (Critical, High, Medium)\\
reason\_statement & String & Formatted, evidence-populated explanation string\\
key\_kpi\_evidence & JSON Object & Dictionary of measured KPI values used in rule matching\\
evaluation\_trace & JSON Object & Complete audit trail of evaluated priority gates\\
recommended\_actions & Array of Strings & Specific engineering steps retrieved from \texttt{solutions.json}\\
standards\_references & Array of Objects & 3GPP and vendor documentation citations\\
\bottomrule
\end{tabularx}
\caption{Suggested downstream integration schema.}
\label[tab]{rca-downstream-schema}
\end{table}
